\chapter{Datos y metodología}\label{cap:datos}

\section{Datos y auditoría}\label{sec:auditoria}

Los datos provienen del \textit{web scraping} de un portal inmobiliario peruano (marzo de 2025): 3\,739 avisos de alquiler de departamentos en 13 distritos, geocodificados con la API de Google. Junto con las coordenadas, la API devuelve su precisión:
\begin{itemize}[itemsep=0pt]
\item \texttt{ROOFTOP}: el edificio exacto;
\item \texttt{RANGE\_INTERPOLATED}: una posición interpolada en la cuadra;
\item \texttt{GEOMETRIC\_CENTER}: el centro de una calle;
\item \texttt{APPROXIMATE}: el centro de la zona.
\end{itemize}
La auditoría (\texttt{notebook2\_auditoria}) aplica, en orden:
\begin{enumerate}[leftmargin=*,itemsep=1pt]
\item exclusión de avisos sin dormitorios, baños, área, precio o distrito;
\item \textbf{deduplicación}: re-publicaciones (mismo texto y atributos; la mitad era el mismo aviso bajo dos distritos) y publicaciones cruzadas (mismas coordenadas, atributos y precio con texto distinto);
\item \textbf{geocodificación incoherente}: más de 8\,km de la mediana de su distrito;
\item \textbf{solo direcciones exactas (ROOFTOP)}, porque la ubicación es el confusor más importante;
\item las reglas de depuración del avance previo (rangos, combinaciones raras, atípicos de precio por m²).
\end{enumerate}
Se crean dos agrupamientos para la inferencia: el \textbf{edificio} (mismas coordenadas) y el \textbf{bloque espacial} de 1\,km. La descripción y la dirección no se exportan.

\section{Medición de la cochera con IA}\label{sec:oe5}

La cochera no es un campo de la base. El procedimiento, implementado en \path{notebook2b_extraccion_ia} y \path{extraccion_ia.py}, minimiza el costo:
\begin{enumerate}[leftmargin=*,itemsep=1pt]
\item una \textbf{regex} busca \emph{cochera, estacionamiento, parqueo, parking, garaje} y sus variantes; sin mención, el aviso se clasifica como \emph{sin cochera} y no se llama a la API;
\item se envía solo el \textbf{fragmento} \texttt{texto[x-200 : x+200]} alrededor de cada mención, con teléfonos y correos enmascarados;
\item el modelo \path{gpt-4.1-mini-2025-04-14}, con temperatura 0, responde un \textbf{JSON de esquema fijo} (\texttt{cochera}, \texttt{tipo}, \texttt{precio\_extra}, \texttt{incluida\_explicita}, \texttt{num\_cocheras}, \texttt{motivo}); un piloto de 25 avisos permitió corregir el prompt;
\item las respuestas se guardan en una \textbf{caché}, de modo que reejecutar no cuesta.
\end{enumerate}
\textbf{Supuestos:} sin mención equivale a sin cochera, y con cochera sin pago adicional equivale a incluida. \textbf{Tratamiento:} $D=1$ si está incluida y $D=0$ si no tiene; se excluyen las de costo extra y las ambiguas. Se preparó una muestra de 249 avisos para validación humana, que está pendiente.

\section{Estimador: Double ML interactivo}\label{sec:irm}

Se usan tres modelos de \textit{histogram gradient boosting} (300 iteraciones, profundidad 4, mínimo de 20 observaciones por hoja), todos con \textit{cross-fitting} de 5 folds \textbf{agrupados por edificio}: cada aviso se predice con modelos que no vieron ni ese aviso ni su edificio.
\begin{itemize}[leftmargin=*,itemsep=1pt]
\item \textbf{Modelo 1, propensión} $\hat m(X)$. Es un clasificador que da a cada departamento su probabilidad de tener cochera. Sirve para descartar a quienes no tienen comparables ($\hat m<0.05$ o $\hat m>0.95$) y para ponderar: un caso poco frecuente en su perfil representa a $1/\hat m$ departamentos.
\item \textbf{Modelos 2 y 3, precio} $\hat g_1(X)$ y $\hat g_0(X)$. Se entrenan con los departamentos con y sin cochera, respectivamente, pero predicen \emph{todos}. Así, cada departamento obtiene su precio estimado con y sin cochera, incluido el contrafactual que no se observa.
\end{itemize}
Para cada aviso en soporte común se calcula el score doblemente robusto, y el efecto es su promedio:
\begin{equation}\label{eq:aipw}
\phi_i=\hat g_1(X_i)-\hat g_0(X_i)+\frac{D_i\,(Y_i-\hat g_1(X_i))}{\hat m(X_i)}-\frac{(1-D_i)\,(Y_i-\hat g_0(X_i))}{1-\hat m(X_i)},
\end{equation}
\[\hat\theta=\frac{1}{n}\sum_i\phi_i.\]
El primer término es la diferencia predicha. Los otros dos corrigen el error del modelo de precio en la versión observada, ponderado por la propensión. Si el modelo de precio o el de propensión es correcto, el estimador es consistente (\emph{doble robustez}). Se repite con 10 particiones y se reporta la mediana. Los errores estándar se agrupan por edificio, y como alternativa por bloque de 1\,km. Los baños se estiman con el modelo parcialmente lineal (dos modelos y regresión de residuos).

\section{Heterogeneidad, sensibilidad y dependencia}\label{sec:sens}

\paragraph{Efectos por grupo.} Los GATE son el promedio de $\phi_i$ por clase de distrito (clase1: Miraflores, San Isidro, Barranco y Lince; clase2: Magdalena, Surquillo, Jesús María, San Borja y Surco; clase3: Cercado, San Miguel, Chorrillos y La Molina) y por quintil de área. La igualdad entre grupos se contrasta con una prueba de Wald.

\paragraph{Sensibilidad.} Para un confusor omitido $U$ con fuerzas $C_Y^2$ (sobre el precio) y $C_D^2$ (sobre el tratamiento), el sesgo queda acotado por \citep{chernozhukov2022}
\begin{equation}\label{eq:cota}
|\hat\theta-\theta|\le\sqrt{\tfrac{C_Y^2\,C_D^2}{1-C_D^2}}\,\sqrt{\sigma^2\nu^2},
\end{equation}
donde $\sigma^2$ es la varianza del error del precio y $\nu^2$ el segundo momento de los pesos. Se reportan el valor de robustez (RV: la fuerza común que anula el efecto) y benchmarks que comparan con la fuerza de los confusores observados.

\paragraph{Dependencia espacial.} Se calcula el I de Moran de los residuos \citep{moran1950}.
